\section{Version history}\label{app:version-history}

This appendix records the differences between successive versions.  No
statement or proof of the paper depends on it.

\subsection{Version 3}\label{subsec:version-3}

Version 3 keeps every numbered result of Version 2, with its name,
number and place.

\paragraph{Results.}
\begin{itemize}
\item \Cref{thm:F13a} takes its source hypothesis \((\mathsf S)\)
  directly as a hypothesis on the interface.  Version~2 stated it for an
  arbitrary source predicate with a certificate and an obligation
  supplying \((\mathsf S)\); the theorem used no property of that
  predicate, so the two forms are equivalent.  The reading paragraph
  gives an interface satisfying \((\mathsf S)\) with a predictive
  surplus.
\item The formal development no longer contains an axiom for
  \Fref{F13a}: the property that Version~2 recorded as a guarded
  trust-base axiom is now the explicit hypothesis \((\mathsf S)\).  The
  opacity-guard subsection and the tagged-axiom register of
  Version~2 are removed accordingly.
\item \Cref{thm:F6} is formalized for real matrices, and over any
  ordered field; a remark after its proof records that the theorem holds
  over any ordered field.  Version~2 formalized it over the rationals
  only, so its formalization is now faithful.
\item The summaries of \Fref{F14} carry the measurability and
  square-integrability of the readout, and the summary of \Fref{F36}
  carries the surjectivity of the limit quotient.
\end{itemize}

\paragraph{Provenance and analogies.}
\begin{itemize}
\item The reading paragraph of \Cref{thm:F14} identifies the odd
  readout of the self-dual trace confinement paper with the case
  \(U=\mathrm{id}\), and distinguishes it from that paper's response
  involution.
\item The description of the Riemann-hypothesis application of
  \Fref{F14} follows the current version of that paper: its conditional
  theorem uses an auxiliary ledger of represented zeros, and is
  equivalent to the Riemann hypothesis on the bare shell of all
  nontrivial zeros; its concrete version uses finite windows of zeros
  closed under \(s\mapsto1-\bar s\), with unit weights, that eventually
  contain every zero, and is likewise equivalent to the Riemann
  hypothesis.
\item The description of the substrate of \Fref{F6} follows the current
  version of the Needle Killer paper, which uses Moore--Penrose
  pseudoinverses of positive semidefinite audit forms.
\item The hidden Markov model analogy of \Fref{F13a} states the
  generic identifiability theorem of Allman, Matias and Rhodes with its
  hypotheses and separates identifiability from estimation; Version~2
  stated a sufficient condition that does not suffice.
\item The Landauer analogy of \Fref{F12} is restricted to the erasure of
  an unbiased bit with a thermal reservoir and states the
  entropy-dependent bound.
\item The matching analogy of \Fref{F9} cites the incentive results of
  Dubins--Freedman and Roth, and the minimal-genome analogy of
  \Fref{F42} no longer asserts the irreducibility of the experimental
  genome.
\item Citations of the companion papers refer to their current versions.
\end{itemize}

\paragraph{Presentation.}
\begin{itemize}
\item The introduction, \Cref{sec:apparatus}, the chapter openings of
  the Access, Stability, Transport and Symmetry chapters, and
  \Cref{sec:anchor:closure-reality-boundary,sec:substrate-instances,sec:sketches,sec:nonclaims}
  are rewritten for readability; \Cref{sec:anchor:closure-reality-boundary}
  is retitled.
\item The \emph{Interaction status} paragraphs, whose first sentence held
  for every law by construction, are replaced by nonclaims paragraphs,
  and \Cref{def:birdint,def:normal-form} and
  \Cref{fig:normal-form-template,fig:axis-topology} are revised
  accordingly.
\item Notes on the formalization are collected in
  \Cref{app:formalization}: the body no longer carries formalization
  footnotes, and the summary tables no longer carry a coverage column.
\item In \Cref{sec:status-records-coherence,sec:self-reference-limits},
  proofs that repeated the descent argument cite the descent criterion
  of \Cref{sec:apparatus}, and the reading paragraphs no longer repeat
  the statements they precede; the split set of \Cref{thm:F19} is named
  \(\Delta_{\mathrm{split}}\).
\item The introduction states that most laws are elementary
  consequences of the descent criterion.
\item The paragraphs on notation and displayed equivalences in
  \Cref{sec:apparatus} are shortened, and the descent criterion is
  stated there with its proof.
\end{itemize}

\subsection{Version 2}\label{subsec:version-2}

Version 2 keeps every numbered result of Version 1, with its name, number
and place. This subsection lists the differences between the two versions:
first the principal change to each result, then further changes to results,
Lean coverage, summary rows and instantiations in the order in which they were
made, and then summaries for the Lean artifact, the apparatus and the
instantiations. The body statements are authoritative; each entry below summarizes a change.

\paragraph{Changes to results.}
\begin{itemize}
\item \Cref{thm:F2} states both repairs by their universal properties: the
  source repair \(\langle q,r\circ F\rangle\) is the coarsest refinement of
  \(q\) through which \(r\circ F\) descends, and a target map \(c\) makes
  \(c\circ r\circ F\) descend exactly when it factors through \(c_F\).
\item \Cref{thm:F3} defines when memory is forced and states the memory
  repair \(\langle\currQ,\predM\rangle\) by its universal property, as an
  instance of \Cref{thm:F2}.
\item \Cref{thm:F4} states its status partition as ``exactly one of four'',
  with each status defined, and states that a globalizing family is
  compatible; its proof gives the exclusivity step that uses this.
\item \Cref{thm:F6} reads its needle clause in the language of
  \Cref{sec:apparatus}. Its proof derives the layer-dissolving decomposition from the projection identity and the monotonicity of the Loewner order under congruence, and the sizes of the bounds and of \(S\) are stated. It defines a needle as a transported dissolving direction exceeding its budget and states that no needle exists.
\item \Cref{thm:F7} defines sufficiency and adds the split-pair form and the
  partition into exactly one of four sufficiency statuses.
\item \Cref{lem:F8} defines retention, definability from the old quotient and
  strictness, assumes the old quotient surjective, proves the split-fiber step
  from non-definability, and adds the partition into exactly one of four
  comparisons. Without surjectivity the equivalence fails: an empty history
  set with a nonempty old quotient and an empty new quotient gives a strict
  extension with no split fiber.
\item \Cref{thm:F9} is stated for the residual records of a package, with the
  inherited rule (RC) and agreement of statuses at a common locus as
  hypotheses, and adds that both hypotheses are sharp.
\item \Cref{thm:F11} names the inherited rule (NO) among its hypotheses,
  states that its factorization is unique, and says which hypothesis gives
  which conclusion. It adds that constancy on the access fibers, emptiness of the overread obstruction and unique factorization are equivalent, and that (NO) holds exactly when every accepted exact current claim factors uniquely through the access quotient.
\item \Cref{thm:F12} defines current visibility, states the properties of the
  financed refinement \(\langle\accessQ,d\rangle\), including its
  minimality, and states its last conclusion under the no-overread rule.
\item \Cref{thm:F13a} names its source predicates in paper terms, adds the
  upstream readout of the determining state, defines quotient exposure,
  predictive collapse and predictive surplus from the maps, and proves its
  fifth conjunct from the fourth; an adopted, guarded source obligation supplies
  the first four conjuncts and surjectivity of the current quotient, and no
  concrete PvNP adapter is constructed. It states the unconditional equivalence of predictive surplus with the failure of predictive collapse for a surjective current quotient, and marks the first four conjuncts as supplied by that obligation.
\item \Cref{thm:F13b} is a direct theorem for a closure with a declared
  determining-state readout: predictive surplus and hidden-state surplus are equivalent, and every surplus pair is separated by the
  determining state. Version~1 stated it under an obligation axiom.
\item \Cref{thm:F14} is a direct theorem: it states the two properties of
  the positive cone and the trace identity its proof uses, with the
  measurability, integrability and equivariance of the readout and proves that a positive anti-invariant
  ledger dominated by bounds of vanishing trace is zero, and that a
  separating readout confines visible mass to the fixed locus. Version~1
  stated it under a source certificate. Its readout is measurable, square-integrable and anti-invariant, \(\psi^-(\Jinv x)=-U\psi^-(x)\), for a declared unitary involution \(U\), and its separation condition is a null-set condition.
\item \Cref{thm:F15a} no longer carries a breaker readout as an unused hypothesis
  (the breaker readout now enters only through the breaker-extension clause), and its
  proof does not use equivariance of the orbit map.
\item \Cref{thm:F15b} lets the formation operator select a branch in
  \(\mathcal C\) itself, with \(q\circ F=q\); in Version~1 it took values in
  the branch quotient, on which the group acts trivially, so the required
  inequality \(g\cdot F(c)\ne F(c)\) could not hold. It is a direct theorem:
  a branch-selection formation record exists exactly when some ground state
  is not invariant under the symmetry, and every such record selects a branch
  in the orbit of the ground state that is not fixed by the group and is not
  the value of any equivariant selector. Its formation record is constant on orbits, so the identity map is not a record, and for a symmetric ground-state predicate the record exhibits two distinct ground states in one orbit.
\item \Cref{thm:F17} states its universal property as ``any quotient that
  factors through every \(q_i\) also factors through \(p\)''; Version~1
  stated the reverse direction, which forces \(p\) to be injective. Part (7)
  of the proof uses descent of \(r\) through \(p\) and of \(p\) through \(r\).
\item \Cref{thm:F18} defines lawfulness of a predicate as invariance under
  the bijections that preserve the access classes, assumes the access quotient
  surjective, and proves that lawfulness is descent.
\item The proof of part (5) of \Cref{thm:F19} uses the obstruction on pairs
  of classes in \(Q_i\times Q_j\). \Cref{thm:F19} declares its target set \(Q_k\).
\item The proof of \Cref{thm:F20} names the constructed descended transport
  \(\theta^{\flat}\) instead of redefining the given map \(\bar\theta\), and
  the theorem adds that, for the descended transport, fact-value stability is
  the raw fact-value equation. Its packaging clause characterizes idempotence by the equality of image and fixed locus.
\item \Cref{thm:F21} states the rules of the audit policy and a conclusion
  about the certifying meta-claim, assumes its three statuses distinct, and
  adds a two-claim model showing that its fourth rule cannot be dropped.
\item \Cref{thm:F23} defines stability as fiber-supported ledgers with the
  declared transport equation, and adds the dichotomy between descent and
  unresolvedness, the value of \(\operatorname{Pr}_{q,r}\) for a descending
  readout, and that the readout distributes the weight of each fiber. It adds that a ledger concentrated on its fiber gives the unrestricted mass of each readout value.
\item \Cref{thm:F22} states the equivalence of an empty mediation obstruction
  with the absence of boundary leakage and of internal insufficiency; Version~1
  stated one direction.
\item \Cref{thm:F24} states that a role split implies the strict layer
  extension \(q_s\); Version~1 asserted that exactly one of eight resolution
  cases holds, and the other cases are now described in the reading paragraph
  as declared statuses. It adds that \(q_s\) is coarsest in the factorization preorder among maps through which \(q\) and \(s\) both factor. Its corestriction to its image is the coarsest surjective quotient with that property.
\item \Cref{thm:F25} defines a causal effect as a stable nonzero contrast
  between two admissible interventions, and adds that, for the descended
  transports, every stable causal effect has raw witnesses.
\item \Cref{thm:F27} defines conservation as invariance under every declared
  step and the orbits as the classes of the generated equivalence.
\item \Cref{thm:F26} defines moduli readouts, presentation dependence,
  necessity and contingency, and proves the descent criterion, the exclusive
  split of a moduli readout into necessary and contingent, and the transport
  criterion; the converse step uses a pair with the same modulus.
\item \Cref{thm:F28} defines composability as descent through parts together
  with interfaces, with associativity and identity as separate conditions,
  and names \(\ell\) and \(\varrho\) as the two bracketings of a triple
  composite.
\item \Cref{thm:F29} defines compatibility of \(\sim\) with \(\pi\), states that
  presentation invariance implies invariance under \(\sim\); when
  \(\sim\) is an equivalence relation, the converse holds for every
  two-valued readout exactly when \(\sim\) is the kernel of \(\pi\).
  It also justifies \(\bar T\circ\pi=T\).
\item \Cref{thm:F30} adds the obstruction criterion for strict compression
  and the continuation clause that uses \(\kappa\), \(e'\) and \(t'\).
\item \Cref{thm:F31} assumes \(q\) and \(\alpha\) surjective, and its
  factorization form requires admissibility to descend; its reading paragraph
  notes that no obstruction uses the target quotient.
\item \Cref{thm:F34} assumes \(\rho\circ\tau\) surjective onto \(R\); in
  Version~1 surjectivity onto its image imposed no condition, and the
  equivalence failed for an empty source with a nonempty recovery quotient.
\item \Cref{thm:F35} assumes the finer scale quotient surjective, which makes
  the descended law unique, and proves both directions of the renormalization
  criterion.
\item The last line of \Cref{thm:F36} states the obstruction criterion for the
  renormalized readout; in Version~1 it was a tautology.
\item \Cref{thm:F37} quantifies over a declared class of admissible joint
  carriers and states separate admissibility; in Version~1 the product map
  was always a joint quotient, so the right-hand side could not hold. It adds
  the split-pair criterion.
\item \Cref{thm:F39} assumes a declared zero of the entropy values, attained
  exactly by quotients with singleton fibers, and adds that, when the mixing
  functional vanishes exactly on pairs with the same fibers, zero target
  entropy is zero mixing \(m(\langle q,t\rangle,q)\), as for conditional
  Shannon entropy. Its mixing clause assumes only that \(m(\langle q,t\rangle,q)\) vanishes exactly when the two quotients have the same fibers, which conditional Shannon entropy satisfies.
\item \Cref{thm:F40} concerns a symmetry that is claimed and not recorded as
  explicitly broken, and adds that the group law holds for the descended
  action of a group action. It defines the anomaly as a nonempty action or domain obstruction or a failure of readout preservation, and the induced action when the action obstruction is empty; with the induced action, an anomaly is exactly the failure of one of the four conditions.
\item \Cref{thm:F41} writes the route actions as \(\mathrm{act}_A\),
  \(\mathrm{act}_B\), states the status-value map between the packages, says
  that role-preserving duality equivalence is defined by the displayed
  conditions, and adds that the inverse map also intertwines the routes. Its role readouts live on the quotients, and the inverse map carries roles back.
\item \Cref{thm:F42} writes the cost of the competing description at \(h\)
  as \(\operatorname{cost}(e'(h))\), with the compression as a subscript of
  the irreducibility predicate, and adds that an irreducible closed
  compression exists when costs are well founded and some compression is
  closed.
\item \Cref{thm:F43} defines the stable limit quotient by its four conditions
  and proves the implication; no limit is constructed.
\item \Cref{thm:F44} writes its defining line with \(:\Longleftrightarrow\),
  states the surjectivity its descent clause uses, and adds that phase
  boundaries and criticality descend when their data do.
\item \Cref{thm:F45} declares the order on local quotients used in its
  minimality clause and defines its locality predicates as factorizations.
\item \Cref{thm:F46} writes its defining line with \(:\Longleftrightarrow\)
  and adds the selected-modulus and transport clauses. It no longer assumes a
  descended readout: descent is fiber constancy, the descended readout is
  unique, and a law-readout is a constant readout.
\item \Cref{thm:F47} writes its defining line with \(:\Longleftrightarrow\)
  and adds the consequences of realized fine-tuning and the monotonicity of
  smallness under a stricter target.
\item \Cref{thm:F49} states the tests for the joint readout together with the
  local access data, adds a declared direct-route explanation, the
  surjectivity hypotheses and a nonempty history space, and classifies an
  unexplained present correlation as a residual, nonlocal exactly when the
  declared locality package fails. Its explanation predicates are defined by
  fiber constancy, so the displayed factorizations are proved equivalences.
\item \Cref{thm:F38} and \Cref{thm:F50} open with the fiber-constancy
  criterion for the transported record and the raw background, and
  \Cref{thm:F48} defines a topological readout as a factorization through the
  gluing quotient and adds that every topological readout is exactly one of a
  complete and a coarse invariant.
\item \Cref{thm:F50} declares its set of formed targets and adds that, for the
  descended background, necessity is constancy of the raw background.
\item \Cref{thm:F51} defines exact unification by the conjunction of
  parent formation, role preservation and closure of the cross residuals, and
  declares the route systems and status translations the display uses; the
  reading paragraph says that exact unification is defined by the displayed
  conditions, and the theorem adds the common-refinement square and the
  inclusion of the parent kernel in both child kernels.
\item \Cref{thm:F52} writes its defining line with \(:\Longleftrightarrow\)
  and adds its obstruction criterion. \Cref{sec:sketches} describes it as a characterization for declared residuals and directions.
\item \Cref{prop:F53} states that its positive equivalence combines the
  closure fixed-point criterion with the definition of the retained operational predicate, with
  the structural defect named. The
  hypotheses of \Cref{thm:CER-direct-no-go,thm:CER-converse-no-go} are
  spelled out before them, the direct one through a declared reduct map, and a
  finite model with a closure operator follows each and satisfies its
  hypothesis. Both theorems state the finite model as their last clause, so
  neither implication follows from the axioms of a closure operator with
  reduct stability, and \Cref{thm:CER-converse-no-go} adds that the converse
  fails exactly when some fixed point of the closure is not bare-realised.
\item \Cref{thm:F10} and \Cref{thm:F5} assume their quotient maps surjective,
  which the coarsest-quotient clause and the recombination criterion need
  when a carrier is empty. \Cref{thm:F13a} states that its certificate is
  used only as the hypothesis of conjuncts 1--4 and that its
  certificate-free content is the last clause. \Cref{thm:F11} and
  \Cref{thm:F9} say which conclusions restate their inherited rules.
  \Cref{thm:F6} states the scope of its substrate Lean anchor, and the
  closure--reality row lists its Lean theorems and says that the positive
  equivalence holds in Lean by unfolding definitions. Summary rows for
  \Fref{F4}, \Fref{F26}, \Fref{F40} and \Fref{F52} match their theorems.
\item \Cref{thm:F13a} describes its source certificate as the guard of an
  adopted source obligation, with no identification with an interface of the
  source paper constructed, and \Cref{thm:F11} is stated
  in three parts: the factorization criterion for every claim, its
  consequence under (NO), and the characterization of (NO). Defining lines of
  \Cref{thm:F29}, \Cref{thm:F36} and \Cref{thm:F40} use
  \(:\Longleftrightarrow\), \Cref{thm:F25} writes its context classes as \(c\),
  and \Cref{thm:F20} and \Cref{thm:F50} name their record-level transport and
  background as declared maps. \Cref{def:diagnosis-source} states the determining hypothesis to which
  \Cref{thm:F13b} applies. \Cref{thm:F15a}
  adds that for the orbit quotient with its induced action, the action on
  orbits is trivial and every equivariant section takes values in \(X^G\);
  when \(O\ne\emptyset\), such a section can exist only if
  \(X^G\ne\emptyset\). The Lean theorems for \Fref{F20}, \Fref{F25} and
  \Fref{F50} assume the descent equation only for the clause that uses it.
\item \Cref{thm:F29} adds section independence, the descent criterion for
  relations on presentation tuples and the independence of equivariance;
  \Cref{thm:F38} defines the arrow at the level of histories and proves its
  record-level form; \Cref{thm:F24} and \Cref{thm:F31} define the strict layer
  extension and counterfactual legitimacy. In \Cref{thm:F36},
  \Cref{thm:F40}, \Cref{thm:F42} and \Cref{thm:F45} definitions precede the
  proved lines. The Lean theorems for \Fref{F11}, \Fref{F23}, \Fref{F24},
  \Fref{F29} and \Fref{F39} carry each hypothesis only where the paper uses
  it. The summary of \Fref{F14} says that the sources take their domination
  bounds as input, and the catalog-boundary definition is linked to
  \Cref{thm:F52}.
\item \Cref{thm:F5} adds the reverse comparison and its four statuses,
  \Cref{thm:F7} states the bundled form of sufficiency, \Cref{thm:F15a} states
  the breaker-extension clauses, and \Cref{thm:F40} lists its four conditions.
  The Lean theorems for \Fref{F7} and \Fref{F17} take the paper's hypotheses:
  the per-probe family and the factorization and minimality of the common
  quotient. \Cref{thm:F13a} says what its certificate identifies and that
  the formal predicate is opaque, the summaries of \Fref{F9} state its
  hypotheses, the displayed-equivalence convention treats every line without
  \(:\Longleftrightarrow\) as an assertion, and Lean names in the body are
  collected in \Cref{app:formalization}.
\item Interpretive sentences that closed the proofs of
  \Cref{thm:F27,thm:F28,thm:F29,thm:F43} follow them.
\item Surjectivity is assumed only where a statement says so
  (\Cref{sec:apparatus}). \Cref{thm:F5,thm:F25,thm:F48,thm:F49,thm:F50} define
  their predicates before the conclusion, and \Cref{thm:F43} defines its test
  set, cofinality and vanishing. The footnotes to \Cref{thm:F13a} and
  \Cref{prop:F53} say which clauses are checked in Lean and which are taken
  from a source paper, the Lean ledger for \Fref{F14} carries the declared
  involution \(U\), and the Access introduction and summary state what
  \Cref{thm:F12} proves. The coverage table follows the chapter order, and
  three instantiations are retagged or reworded to match their content.
\item \Cref{thm:F50} classifies the vacuum criterion as empty, singleton or
  multiple and states existence and uniqueness of a vacuum;
  \Cref{thm:F25} defines its descent predicates and gives each raw-witness
  clause only its own descent equation; \Cref{thm:F38} assumes surjectivity of
  the source record quotient only. \Cref{def:closure-formation} defines the
  formed closure of a rule set. The Lean theorems for \Fref{F9},
  \Fref{F25}, \Fref{F38} and \Fref{F50} state these forms, with an arbitrary
  status type for \Fref{F9}; the coverage row of \Fref{F53} lists first the
  theorems that take the paper's hypotheses. The formalization appendix
  states that the trust-base axiom can refute its guard, the answer to the
  anchor-row objection says what each anchor row adds, and the catalog
  application of \Cref{thm:F52} states which of its clauses are met.
\item \Cref{thm:F2} states the factorization of the source repair through
  every refinement, \Cref{thm:F45} adds that the local repair carries the
  target, and \Cref{thm:F51} derives its square and kernel inclusions from
  the projection equations alone. \Cref{thm:F14} states its conclusion up to
  null sets, and the readings of \Cref{thm:F20,thm:F23} say what their
  predicates join. The Lean theorems for \Fref{F2}, \Fref{F9}, \Fref{F12},
  \Fref{F36}, \Fref{F43}, \Fref{F44}, \Fref{F45} and \Fref{F51} state every
  displayed clause, \Fref{F43} with index-dependent stage and residual
  types; the coverage row of \Fref{F53} lists only applicable theorems, and
  the formalization appendix states what the opaque guard of \Fref{F13a}
  permits.
\item \Cref{thm:F13a} lists surjectivity of the current quotient among the
  imported data, \Cref{thm:F33} states the family clause as an equivalence
  for the bundled readout, and \Cref{thm:F6,thm:F37} say which setting words
  add no condition. The abstract and introduction count the laws proved in
  full and the laws with direct Lean theorems separately. The Lean
  development instantiates the \Fref{F14} laws for positive semidefinite
  rational matrices, proves the \Fref{F33} family clause, and proves the
  fixed-point form of the \Fref{F53} criterion for the structural defect of
  the retained five-flag encoding. The coverage vocabulary is Theorem, Anchor and Partial
  with alignments faithful, restricted and anchored, and the build
  instructions are replaced by a statement of the trust base.
\item The readings of \Cref{thm:F17,thm:F19,thm:F20} name the case of each
  status, the comparisons with G\"odel's and Noether's theorems and with the
  information paradox state what each classical result derives, the framing
  of \Cref{thm:F37} states that non-joint access is relative to the declared
  class of joint carriers, and the prose around \Cref{thm:F6} states what the
  theorem proves. The Lean theorem for \Fref{F52} allows an arbitrary verdict
  set, and the coverage rows of \Fref{F8} and \Fref{F42} list the theorems
  that carry their hypothesis-free clauses.
\item \Cref{thm:F49} states when a correlation residual is present including
  the admissibility and stability of the route, the \Fref{F21} nonclaim names
  the rules its conclusion uses, and the reading of \Cref{thm:F19} gives its
  obstruction sets in ordinary notation. Surjectivity is assumed only where a
  statement says so, a faithful Lean row may state more than the paper, and
  the build statement names its scope. The abstract counts the laws in one
  split, and the presentation figure describes where each part of a law
  appears.
\item \Cref{thm:F35} states the next-scale descent criterion that its
  linking equation gives, \Cref{thm:F23} characterizes stable probability as
  stability together with failure of descent, \Cref{thm:F40} names the
  induced action separately from the candidate action, and the residual of
  the closure--reality apparatus has a stated codomain. Two instantiations
  whose data do not meet their law's hypotheses are tagged as analogies, and
  the abstract states the scope of the formalized matrix inequality.
\item The Myhill--Nerode, crystallographic, Gr\"obner-basis and general
  relativity instantiations are restated so that they meet the hypotheses of
  their laws; \Cref{thm:F47} keeps only the predicates its conclusion uses;
  gates are described as maps on carriers; the introduction describes the
  eight themes as chapter labels; the abstract counts the laws proved in full
  separately from those with direct Lean theorems; and two displays are set
  within the margins.
\item \Cref{thm:F42} measures irreducibility against a declared class of
  admissible compressions and states what it reduces to when every map is
  admissible, and \Cref{thm:F21} shows that rule~3 cannot be dropped either;
  the Lean development proves both. The residual of the closure--reality
  apparatus is the set of its five structural defect atoms, the hiddenness
  interface defines \(\operatorname{HiddenFromCurrent}\) with its other
  predicates, and a paragraph on the pairing refinement precedes the
  catalog. The readings of \Cref{thm:F24,thm:F46} relate them to
  \Cref{thm:F12,thm:F26}, the similarity instantiation of \Cref{thm:F29}
  states that its listed invariants do not separate similarity classes, and
  the build statement names Lean's standard axioms.
\item \Cref{thm:F23} adds a clause for finite weights, in which descent is
  equivalent to every fiber probability being a point mass, and its
  concentration clause is proved from additivity alone; \Cref{thm:F43}
  characterizes the stable limit quotient for every index;
  \Cref{thm:F45} characterizes minimal domains under strict factorization;
  \Cref{thm:F51} states when the parent carries exactly the children's
  distinctions; and \Cref{thm:F52} shows that each of its four conditions is
  independent of the others. The Lean development proves these clauses. The
  catalog boundary, horizon-bounded support, record formation and diagnosis
  source data are defined by formulas, \Cref{thm:F14} assumes only that the
  traces have infimum zero, and the quantum instantiation of
  \Cref{thm:F19} places decoherence in the dissolution obstruction.
\item \Cref{thm:F27} allows partial steps, so that moves applying only where
  they apply are instances, and \Cref{thm:CER-converse-no-go} shows that both
  closure--reality implications fail for every nontrivial closure operator on
  subsets; the Lean development proves both. The Reed--Solomon instantiation
  places the error vector in the source carrier, the introduction says how the
  eight themes are assigned to chapters, the coverage row of \Fref{F23} names
  the theorem for each clause, the formalization appendix gives each restricted
  or anchored row its own paragraph, and the pairing paragraph lists every law
  that uses the pairing.
\item \Cref{thm:F23} defines its transported ledger and target mixture
  through a ledger transport, a target mixture and the class-level
  continuation, \Cref{thm:F49} defines its correlation record from the joint
  readout and the local access classes, and the duality ledger of
  \Cref{thm:F14} is strongly measurable. The crystallographic, Burnside and
  quantum instantiations state the predicate or functional they use, the
  abstract states the rational-matrix scope of \Fref{F6}, the formalization
  appendix distinguishes the inherited diagonal-support proof from the
  Foundations IV endpoint proof for arbitrary rational blocks, the \Fref{F13a} note names its
  axiom, guard and axiom-free clause, the introduction relates its coverage
  counts to those of the abstract, and two summary-table rows state what their
  theorems prove.
\item The closure--reality anchor defines the retained operational predicate by the vanishing
  of the residual in the body and in the proof, \Cref{thm:F20} says when its
  three stable-fact conditions concern one recorded fact, and
  \Cref{thm:F33} states horizon-boundedness and defines horizon
  sufficiency. The inherited-rules paragraph counts the laws that use them,
  the source of \Fref{F13b} is called the CSL paper, access is stated for any
  observation signature, the Access chapter points new readers to the
  Transport chapter, and three summary-table rows state what their theorems
  prove.
\item The Lehmann--Scheff\'e instantiation uses the likelihood-ratio function,
  whose fibers are those of the minimal sufficient statistic, and the Markov
  instantiation uses the transition kernel of a countable chain. The
  closure--reality chapter states what the no-go theorems show and that
  the retained operational predicate is not compared with bare realisability, the
  formalization appendix states the restriction of the two anchored rows,
  the abstract's coverage sentences are regrouped, the chapter groupings are
  called axes throughout, the covering condition of \Cref{thm:F43} is
  explained, and the predictive-state and residual-coverage statements are
  corrected.
\item \Cref{thm:F44} shows that a phase transition along a path of
  neighbouring formed controls crosses a phase boundary, \Cref{thm:F39} adds
  the fiber-volume instance, and \Cref{thm:F13a} states the exposure,
  collapse and surplus equivalences under the certificate; the Lean
  development proves these clauses, proves the \Fref{F6} endpoint for
  arbitrary rational blocks, and defines closedness in the closure--reality
  anchor as a fixed point of the saturation operator. The \Fref{F13b}
  statement says that surplus pairs are separated by the determining state,
  the proof of \Cref{thm:F11} cites the no-overread rule, the reading path
  names the one cross-chapter use, and each Lean footnote names its row of
  the coverage table.
\item \Cref{thm:F44} adds the clause for continuous paths in a topological
  control space, and \Cref{thm:F42} states its collapse clause for any
  admissible class that isolates \(h\) with every value; the Lean
  development proves both. \(\BirdInt\)-compatibility is defined by listing in
  the law's setup, \Cref{thm:F40} declares the record of explicit breaking,
  \Cref{thm:F15a} assumes the orbit-quotient property only where it is used,
  the statements of \Cref{thm:F17,thm:F19,thm:F36} gloss their predicates,
  the reading paragraph of \Cref{thm:F43} names every conjunct unused by the
  construction, and the summary rows of \Fref{F15a} and \Fref{F39} follow
  their statements.
\item \Cref{prop:F53} adds its clause on the saturation of the retained five-flag encoding,
  and the closure--reality chapter states the general content of its two no-go
  theorems. The notation of \Cref{thm:F6} allows any symmetric \(C^{+}\) and any
  symmetric \(K_{LL}^{\dagger}\) satisfying the projection identity, the reading
  paragraphs of \Cref{thm:F13a,thm:F39,thm:F42} state what their predicates do
  and do not require, the introduction and abstract state which laws import
  results and how Lean covers them, the Access chapter can be read first, and
  the interaction-status convention adds no premise.
\item \Cref{thm:F13b} characterizes determining readouts by factorization
  and splits predictive surplus for an arbitrary readout, \Cref{thm:F29}
  characterizes presentation invariance by a generating relation of the
  kernel, and \Cref{thm:F36} relates the completion of a renormalized target
  obtained by post-processing to that of its raw target; the Lean development
  proves the three clauses. The encoding of the reflective closure used by the
  saturation clause is described with the definition, every
  interaction-status paragraph names its setup data, the instantiations of
  \Fref{F46} state that the classical law is an input, and the reading
  paragraphs of \Cref{thm:F43,thm:F44} state what their conclusions add.
\item \Cref{thm:F42} adds its global form over a base description,
  \Cref{thm:F19} shows that holonomy is coherent when the three class maps are
  the descended ones, \Cref{thm:F14} requires only that the traces approach
  zero from above, the no-overread rule names its acceptance predicate, the
  fiber-volume instance of \Cref{thm:F39} states where its mixing term is
  evaluated, the formalization appendix names the declarations of
  \Fref{F13a}, and every law is noted as \(\BirdInt\)-compatible by
  construction.
\item \Cref{thm:F42} adds the existence of a globally irreducible description
  of minimal cost, \Cref{thm:F47} adds a quantitative clause for
  complementary selector regions, \Cref{thm:F39} requires combine-monotonicity
  only at the pair it uses, so that the fiber-volume instance satisfies the
  theorem, the closure--reality no-go theorems hold for every extensive
  operator, and the Lean development proves these clauses and the full
  descent equivalence of \Cref{thm:F35}. \Cref{thm:F43} states its unique
  compatible readout as the conclusion, the direct no-go theorem characterizes
  failure of its implication, the vocabulary defines a closure space, the
  reading path lists the arguments that are not descent arguments, and the
  Reflexive Nonclosure law states which of its conclusions are premises.
\item The point-mass clause of \Cref{thm:F23} requires weights only on each
  fiber, so distinct classes may share a ledger, and the Lean development proves
  it in that form. \Cref{thm:F6} is stated for arbitrary block matrices with the
  symmetry and projection identities it uses, \Cref{prop:F53} displays the
  fixed-point equivalence it asserts, two instantiations are tagged as
  analogies, and \Cref{thm:F17} writes its defects out.
\item \Cref{thm:F13b} identifies the pairing with the predictive quotient as
  the coarsest determining one, \Cref{thm:F49} shows that a correlation
  residual is exactly the absence of a local explanation, and \Cref{thm:F41}
  shows that the inverse class map and the role-value map are determined; the
  Lean development proves these clauses. The definition of the current and
  predictive quotients states what a quotient and refinement are, the
  interaction-status convention states that it adds no hypothesis, the
  closure--reality proof states what the no-go theorems show, and the coverage
  rows of \Fref{F23} and \Fref{F39} name the declarations of their clauses.
\item The closure--reality no-go theorems characterize the extensive operators
  for which a reduct-stable predicate makes both implications fail, and the Lean
  development proves the characterization. \Cref{thm:F18} states its invariance
  condition in the statement, \Cref{thm:F39} defines its entropy predicates in
  the statement, \Cref{thm:F44} states what its boundary means for continuous
  paths, the settings of \Cref{thm:F9,thm:F11,thm:F41} are stated as intended
  readings, and the nonclaims name the \(\BirdInt\) judgment as background.
\item \Cref{thm:F44} asks along a finite path only that each control have
  neighbours of the next control's phase, which in the topological case means
  that it lies in the closure of that phase class, and \Cref{prop:F53} shows
  that both closure--reality implications hold exactly for the fixed-point
  predicate of the closure, so that the retained operational predicate is the
  unique predicate for which closure equals reality holds for
  \(\operatorname{Sat}_5\); the Lean development
  proves both. The reading of \Cref{thm:F17} is stated in ordinary terms, the
  footnote to \Cref{thm:F14} describes its Lean instance exactly, and the
  Vitali and sheaf-locality instantiations name the readout and the minimal
  domain.
\item The introduction identifies the inherited no-overread rule (NO)
  as a premise in \Fref{F11} and \Fref{F12} and summarizes the coverage
  in two sentences; the direct
  closure--reality no-go names the carrier at which the implication fails;
  clause (d) of \Cref{prop:F53} is stated for every operator; the reading of
  \Cref{thm:F44} explains its neighbourhood relation; \Cref{thm:F43} glosses
  its covering condition in the statement; and the summary rows of
  \Fref{F8}, \Fref{F23} and \Fref{F44} state the results rather than their
  definitions.
\item The Lean development proves the characterization of the phase boundary
  in clause (iv) of \Cref{thm:F44} and the closure form of the step hypothesis
  of clause (iii). \Cref{thm:F42} defines its obstruction sets in the
  statement, \Cref{thm:F27} states its route structure in the statement, the
  reading of \Cref{thm:F49} explains its surjectivity hypothesis, the anomaly
  clause of \Cref{thm:F40} is stated for unbroken symmetries, the interaction-status
  convention says that the paragraphs may be skipped, and the direct
  closure--reality no-go states that its first clause holds by definition.
\item The Lean development proves the two factorizations of the canonical
  pairing in \Cref{thm:F13b}. The reading of \Cref{thm:F15b} states when a
  non-invariant ground state is a spontaneous breaking, \Cref{thm:F26} explains
  its obstructed clause, \Cref{thm:F19} gains a summary in ordinary terms, the
  Access chapter cites the apparatus pairing refinement in place of
  \Cref{thm:F3}, and the reading groups
  of the Status, Records and Coherence chapter are listed.
\item The reading of \Cref{thm:F23} states which data its stability condition
  constrains; \Cref{prop:F53} describes its right-hand side as the defect
  characterization of the closure's fixed points; the pointwise clause of
  \Cref{thm:F42} states what it depends on for large admissible classes; the
  opening of the Self-Reference and Limits chapter states the horizon law as
  proved and locates the self-audit law; the certificate of \Cref{thm:F13a} is
  described as fixing one interface; and the introduction explains the
  F-codes.
\item \Cref{prop:F53} adds clause (e): both closure--reality implications fail
  for the retained saturation \(\operatorname{Sat}_5\) itself, for a
  reduct-stable predicate; \Cref{thm:F37}
  states its factorization form and the product corollary; and the Lean
  development proves these and the empty-moduli clause of \Cref{thm:F26}. The
  reading of \Cref{thm:F40} states what the recorded flag affects, the
  correlation predicate of \Cref{thm:F49} is described as defined from the
  data, and one instantiation of the Self-Reference and Limits chapter is
  relabelled an analogy.
\item The Lean footnotes point to the per-row map of the formalization
  appendix; the dependency hints of the axis figure state the proof-level uses;
  the convention on displayed lines names the lines that unfold a definition;
  \Cref{thm:F28} defines its associativity and identity predicates; the
  comparison statuses of \Cref{thm:F19} are stated with the theorem; the
  reading of \Cref{thm:F23} states which clauses use the continuation data; and
  the coverage row of \Fref{F53} lists the finite instances.
\item \Cref{thm:F4} adds the characterization of the sheaf condition for a
  cover, and \Cref{thm:F14} adds, for \(U=\mathrm{id}\), the equivalence between
  a vanishing ledger and full measure of the fixed locus; the Lean development
  proves both (the second for the abstract ledger under explicit laws). The
  converse closure--reality no-go identifies its closure-instantiation-gap
  hypothesis with a witnessing carrier, \Cref{thm:F23} names the operation and zero of clause (v),
  the global clause of \Cref{thm:F42} drops an unused surjectivity, and the
  reading order states that the Transport chapter can be read first.
\item Clause (e) of \Cref{prop:F53} names its predicate and requires only a
  carrier over the scope; \Cref{thm:F19} defines its persistence predicates in
  the statement; the last line of \Cref{thm:F22} holds for every \(\bar U\); the
  reading paragraphs of \Cref{thm:F15a,thm:F44} name their symbols in order;
  the anchor status describes what the Lean proves for \Fref{F6} and
  \Fref{F53}; the interaction-status convention points to the nonclaims it
  carries; and the reading of \Cref{thm:F49} states when the source map is
  injective.
\item The convention on displayed lines adds examples of definitions and of
  lines combining definitions with descent criteria; \Cref{prop:F53} states at its head which clauses are imported,
  definitional and proved; the interaction-status paragraphs of the Access,
  Stability and Symmetry chapters are titled with their nonclaims; the proof
  of \Cref{thm:F6} calls its identity the projection identity; and the general
  clause of \Cref{thm:F43} names its hypotheses.
\item The special-case readings of the catalog rows (query determinacy, the
  gas and Coulomb models, connected genus, the Riemann and removable-limit
  quotients, the pyramid carrier) state the hypotheses under which they
  instantiate their rows; \Cref{thm:F43} separates its hypotheses from its
  conclusion; and \Cref{thm:F25} states the uniqueness of the descended
  outcome maps, with the full Lean statement cited in its coverage row.
\item The final clause of \Cref{thm:F14} holds whenever \(U\) has no
  \(-1\)-eigenvector, with a one-point ledger showing the condition is needed;
  \Cref{prop:F53} states the uniqueness of the retained predicate in clause (c)
  and defines the flag set and reducts of clause (e); the proof of
  \Cref{thm:CER-direct-no-go} points to where its extensive-operator clause is
  proved; the substrate table describes what \Fref{F53} imports; and the
  Berry-phase, cryptographic-hash and Voronoi special cases and the weight
  hypothesis of \Cref{thm:F23}(vi) are stated exactly.
\item The obstructed-closure part of \Cref{thm:F26}(iii) holds without the
  standing assumption, which then forces the closure space to be empty;
  \Cref{thm:F20} adds that a stable fact transports the event it records;
  the closure--reality chapter says which operator ``retained'' and
  ``current'' refer to; and the meiosis, abstract-data-type, de Bruijn,
  Vitali, martingale and fundamental-group special cases are corrected.
\item Special cases state the data of their laws and the hypotheses they
  need: the elliptic boundary-value problem (bounded Lipschitz domain,
  coefficient and source classes, Neumann compatibility), the holonomy,
  Berry-phase and Preisach cases of \Fref{F3}, integrability of the
  martingale, the minimal domains of \Fref{F45}, the Bayes closure space,
  the Euler-characteristic transport and the path-integral readout.
\item \Cref{thm:F46} characterizes state-dependent readouts by pairs of
  distinct moduli classes; \Cref{thm:F41} adds that the dualities preserve
  quotient fibers; \Cref{thm:F22} states uniqueness of the mediated update;
  \Cref{thm:F37} drops the surjectivity of its access quotients;
  \Cref{thm:F39} states the laws its fiber-volume instance satisfies; the
  terms of clause (b) of \Cref{prop:F53} are defined; and the proofs of
  \Cref{thm:F15a}, \Cref{thm:F17}, \Cref{thm:F31}, \Cref{thm:F42} and
  \Cref{thm:F44} state only what they prove.
\end{itemize}

\paragraph{Lean artifact and coverage.}
The catalog has 49 direct theorem rows, two anchor rows, and one conditional
row. Each direct row has a Lean theorem, and 48 direct rows have faithful Lean
coverage of the displayed result; \Cref{app:formalization} gives the alignment
of every row. Version~1 described all direct rows as faithful and
listed the declarations for \Fref{F13b} and \Fref{F15b} as theorems, although
they were obligation axioms. Version~2 removes both obligation axioms, proves
\Fref{F13b}, \Fref{F14} and \Fref{F15b} without a guarded axiom, and adds Lean
theorems for the clauses added above. The Lean theorems for \Fref{F11},
\Fref{F17} and \Fref{F21} take the paper's hypotheses and rules explicitly
instead of reading the conclusions off record fields, and the closure--reality
no-go theorems are proved over an abstract carrier and instantiated by the
finite model. One conditional law (\Cref{thm:F13a})
and one trust-base axiom remain; the axiom supplies only the four
source-specific conjuncts of \Cref{thm:F13a} and surjectivity of its current
quotient.

\paragraph{Apparatus and instantiations.} The current and predictive quotient maps are written \(\currQ\) and \(\predM\), distinct from their target sets, and defining lines in several laws use \(:\Longleftrightarrow\). Three instantiations whose data are
empirical are tagged as analogies, and the meiosis instantiation counts
crossovers by parity.
\Cref{sec:apparatus} gives a reading guide: the chapters and axes, the status
vocabulary, the acronyms, the meaning of the layer instantiations, a reading
path and a map of results with the section of each result. It adds paragraphs
on the conditional laws, the two rules inherited from earlier papers,
displayed as formulas, the words that name a law's setting, the emptiness
predicates, obstruction sets, reused letters and displayed equivalences. A
small running example is carried through the introduction and the chapter
openings. A paragraph \emph{Reading the statement} precedes each result whose
statement uses predicates that need explanation, and a paragraph \emph{What the
source supplies} follows each conditional law. The chapter titles name their
subject, and the conditional laws are marked in their subsection titles. The
summary rows paraphrase the statements, symbols that named two objects were
renamed, and the result codes link to their results. The interaction status
and nonclaims of each row carry run-in labels. Each layer instantiation is
tagged as a special case or an analogy; seventy-three instantiations were
corrected where they misstated a condition, a formula, a direction, a date or
an attribution, and two references were added. Each law keeps its special cases beside it, and every law with analogies keeps at least one beside it (\Fref{F29}, \Fref{F45} and \Fref{F46} have special cases only); the other analogies are collected in
\Cref{app:further-instantiations}. The substrate map is given once, in
\Cref{tab:substrate-instances}.

\paragraph{Development of the Lean artifact.}
Before Version 1, an axis-level review of the Lean development led to three
strengthened statements. For \Fref{F2}, the canonical target repair is built
from the equivalence generated by identifying \(rF(x)\) and \(rF(x')\)
whenever \(q(x)=q(x')\), matching the minimal-repair universal property of
the statement, in place of a fallback to a one-point target. For \Fref{F10},
the access quotient \(q_I\) is constructed from a lawful observation
signature rather than assumed, and the universal property covers observable
factorization together with the coarsest-quotient property. For \Fref{F53},
the two closure--reality no-go statements are derived from guarded
structural assumptions (reduct stability on the direct side and a closure
instantiation gap on the converse side) rather than assumed as raw
counterexamples.

\paragraph{Closure operators of the closure--reality chapter.}
The current revision distinguishes the retained five-flag saturation
\(\operatorname{Sat}_5\) from AOR's current evidence-preserving inference
reflector. Earlier descriptions identifying those operators are superseded
by \Cref{def:reflective-closure,prop:F53}. The retained saturation theorem
and CER controls remain valid. Current AOR full accounting has its separate
source-witness criterion, and a Lean counterexample establishes that inference
closure alone does not pay the obligation \(5\le2\).

\paragraph{Source scope of F13a and matrix scope of F6.}
\Fref{F13a} remains conditional on an adopted guarded source obligation.
The current PvNP Lean does not construct the required catalogue adapter
or its opaque certificate; earlier descriptions of its four clauses as
proved source imports are superseded by
\Cref{thm:F13a,subsec:narrowing-disclosure}.
For \Fref{F6}, arbitrary symmetric blocks support the algebraic budget
theorem, while the optimization and positive-residual interpretation
requires the positive-semidefinite setting \(C^{+}\succeq0\) of the
Notation paragraph before \Cref{thm:F6}.
